\slajdid{09_3-s005}
\begin{frame}[label=09_3-s005]
\gusttekst\setlength{\abovedisplayskip}{3pt}\setlength{\belowdisplayskip}{3pt}\setlength{\abovedisplayshortskip}{2pt}\setlength{\belowdisplayshortskip}{2pt}\textbullet\quad Zbog ovih kapacitivnosti pojavljuju nam se \alert{DVA} problema.
\par\medskip
\alert{1.} Koliki treba da je odnos $k_p$ i $k_n$ (odnosno $W_p$ i $W_n$) \underline{u jednom} CMOS invertoru da bi se minimiziralo kašnjenje kada taj invertor treba da puni i prazni spoljne kapacitivnosti, koje su posledica ulaza u naredna logička kola.
\par\smallskip\centering\begin{tikzpicture}[x=.8cm,y=.65cm,font=\sitneoznake,>=Latex]
\ctikzset{bipoles/length=.6cm}
\foreach \x in{0,2.2}{\draw(\x,-.4)--(\x+1,.2)--(\x,.8)--cycle;\draw(\x+1.1,.2)circle(.1);}
\draw(-.6,.2)node[left]{$v_I(t)$}--(0,.2);\draw(1.2,.2)--(2.2,.2);\draw(3.4,.2)--(4,.2)node[right]{$v_O(t)$};
\draw(2.4,-.28)to[C,l=$\mathrm{Cint}$](2.4,-1.35)node[ground]{};
\end{tikzpicture}
\par\medskip\raggedright
\alert{2.} Koliki treba da je međusobni odnos dimenzija tranzistora \underline{među} CMOS invertorima (unutar jednog ćemo odrediti odnose $W_p$ i $W_n$ rešavajući prethodni slučaj) da bi se minimiziralo kašnjenje \underline{u lancu} invertora koji treba (lanac treba) da puni i prazni spoljne kapacitivnosti, a da ulazna kapacitivnost u lanac ostane minimalna.
\par\smallskip\centering\begin{tikzpicture}[x=.8cm,y=.65cm,font=\sitneoznake,>=Latex]
\ctikzset{bipoles/length=.6cm}
\foreach \x/\n in{0/1,1.7/2,5/n}{\draw(\x,-.4)--(\x+1,.2)--(\x,.8)--cycle;\draw(\x+1.1,.2)circle(.1);\node at(\x+.3,.2){$\n$};}
\draw(-.65,.2)node[left]{$v_I(t)$}--(0,.2);\draw(1.2,.2)--(1.7,.2);\draw(2.9,.2)--(3.5,.2);\node at(4.1,.2){$\cdots$};\draw(4.6,.2)--(5,.2);\draw(6.2,.2)--(7.5,.2)node[right]{$v_O(t)$};\draw(6.9,.2)to[C,l=$C_L$](6.9,-1.2)node[ground]{};
\node[below]at(.35,-.6){1};\node[below]at(2.1,-.6){$1\times f$};\node[below]at(5.4,-.6){$1\times f^{n-1}$};
\draw[->](-.85,-1.25)node[below]{$C_{\min}$}--(-.85,-.1)--(-.4,-.1);
\end{tikzpicture}
\par
$f$ – faktor povećanja istovremeno i p i n tranzistora u narednom invertoru
\end{frame}
